\section{Conclusion}
\label{sec:conclusion}

Three empirical findings constitute the core contribution of this paper. First, creative
quality in both authoring and reviewing roles is governed by a three-tier cognitive
architecture: orientation (Tier 1, irreplaceable), operational discipline (Tier 2,
load-bearing), and domain specificity (Tier 3, amplifying). These tiers are not
additively independent --- Tier 3 content without Tier 1 activation degrades performance
below the no-context baseline, and Tier 2 components without Tier 1 direction operate on
the wrong objective. Second, domain expertise is Tier 3 in both roles. It is necessary
for achieving ceiling performance, substitutable across sources and formulations, and
never sufficient without the orientation and discipline layers that determine what the
domain knowledge is being used for. Third, the minimum sufficient creative prompt
consists of three sentences: ``Think first about how the other person will experience
this. Remove anything unnecessary. Verify every claim.'' This three-sentence OLE formula
--- Orientation, Lens, Expertise --- crosses the publishable quality threshold in
the authoring study and is derivable from the reviewing-side tier structure independently.

The practical implication is direct. Any practitioner using an AI system for creative
work can restructure their prompts to follow OLE order: orientation sentence first,
discipline sentences second, domain content as context third --- never as persona identity.
The evidence base for this restructuring is quantitative and large. The orientation
sentence alone accounts for a 15.6-point quality gap between the artist and puzzle
designer personas on a 45-point scale. Its absence accounts for a 29.4-point catastrophic
failure that produces categorical breakdown rather than degraded output. The minimum
sufficient set (T1 + T4 + T5) reaches the publishable threshold at 35.7 out of 45 in the
authoring study, and the equivalent structure (Tier 1 + three non-redundant Tier 2
principles) reaches the best-condition performance in the reviewing study. The OLE
formula is not a heuristic extracted from subjective experience; it is a description of
the empirically measured quality hierarchy across two roles in a controlled testbed.

Three directions for future work extend the tier model beyond puzzle-hunt design. First,
cross-domain validation: does the OLE reconstruction curve replicate in business document
generation, legal drafting, or interactive fiction? The tier model makes a strong
prediction --- orientation-first prompts will outperform domain-expertise-first prompts
in any creative domain whose quality criterion is defined by receiver experience --- and
this prediction is falsifiable by any controlled authoring study that varies prompt
structure rather than prompt content. The tier model's domain-independence is a structural
prediction grounded in the universal role of creative receivers, not an empirically
established finding; cross-domain replication is the required next step.
Second, multi-agent creative systems: when a
reviewing agent and an authoring agent both operate from the OLE structure, does the
quality of their collaborative output exceed that of domain-labeled agents? The symmetric
architecture of authoring and reviewing tiers predicts a compounding effect: an authoring
agent oriented toward the solver and a reviewing agent oriented toward the solver and
author produce a feedback loop in which Tier 1 orientation is reinforced at every
exchange. Third, human validation: do human professionals who explicitly apply the OLE
structure to their own creative briefs report quality improvements relative to their
domain-expertise-first defaults? The tier model predicts yes in all three cases, and in
all three cases the prediction is specific enough to be falsified. A tier model that
fails to cross the threshold at the Tier 1 + Tier 2 boundary in a new domain, or fails
to produce quality advantages in a multi-agent OLE system, would require revision. The
value of the model is precisely its falsifiability --- and the precision with which it
specifies where and how domain expertise should enter a creative prompt.
